\section{Problems: the same equations in both codes}
\label{sec:problems}

All four problems are multigroup diffusion on homogenized regions, and each
is posed to both codes with identical data. Discretization is the only
difference. \ffn{} uses continuous-Galerkin finite elements of degree $p$ on
hexahedra. \garnet{} uses second-order cell-centred finite volumes on
Cartesian boxes or triangular prisms. Both therefore converge to the same
continuous solution.

\paragraph{IAEA-3D.} The two-group benchmark (ANL-7416 Suppl.~2) on a
$17 \times 17 \times 19$ lattice of 20~cm nodes. \ffn{} runs its own input
(\code{examples/3D\_IAEA}). That input applies the albedo condition
$D\,\partial\phi/\partial n = -0.4695\,\phi$ on every face of the full core,
and the corners of the octagonal core are holes. \garnet's
\code{build\_iaea} carries the same data. Its albedo is the same Robin
coefficient: \ffn{} \code{matrix\_operators\_petsc.cc} adds
$\alpha \int \phi\,v$ on boundary faces, and \garnet{} passes $\alpha$ to its
Robin face term. The published $k = 1.02903$ assumes a different boundary
condition. It is quoted for context only and is not used as a reference.

\paragraph{3-D Langenbuch at $t = 0$.} The two-group LMW core
($11 \times 11 \times 10$ nodes of 20~cm) with zero-flux boundaries. Bank~1
is half inserted, with its tips at 100~cm on a node face, and bank~2 is
withdrawn, so no node is partly rodded. \ffn{} runs
\code{examples/3D\_Langenbuch} with its \code{.bar} file at the initial
positions. Two discrepancies in \garnet's builder were found and removed
before any timing:
\begin{itemize}
\item \ffn's map excises the four corner columns. \garnet{} had filled them
  with reflector. \code{build\_langenbuch(excise\_corners=True)} now poses the
  staircase zero-flux boundary instead; the default is unchanged.
\item \garnet's cross sections had been transcribed with one digit fewer than
  the \code{.xsec} file (a relative difference of about $10^{-8}$). They now
  carry the file's digits.
\end{itemize}

\paragraph{Input identity, IAEA-3D and Langenbuch.}
\code{tests/validation/test\_femffusion\_equivalence.py} rebuilds \ffn's node
map from the \code{.prm} \code{Geometry\_Points} and the \code{.xsec}
\code{PLANE} blocks, and applies the bars at their initial heights. It then
checks node by node that \garnet's material map, active mask, cross sections
and boundary conditions are the same. \garnet's first array axis is \ffn's
$y$, so \garnet's maps are \ffn's transposed in plan. The two models are
mirror images across the core diagonal: $k$ is identical, and node powers are
compared after the transpose.

\paragraph{HP-MR, hex-homogenized (11 groups).} The ANL/INL heat-pipe
microreactor: 30 fuel hexagons, a central cell, 12 beryllium reflector
hexagons and 12 control drums, on a 26.752~cm lattice. The core is 200~cm
high, with 20~cm beryllium axial reflectors. Cross sections are the 11-group
ENDF/B-VIII.0 library of the INL Virtual Test Bed, with the fuel assembly
flat-flux homogenized. Each hexagon is one material in each axial zone. The
drums are withdrawn: a drum hexagon is the library's drum beryllium, with no
B$_4$C (Section~\ref{sec:blackdrum} explains why). The boundary is
vacuum (Marshak, $\alpha = 1/2$) on every face in both codes. \ffn's input is
written by \code{tools/export\_femffusion.py}:
\begin{itemize}
\item a Gmsh mesh of three quadrilaterals per hexagon, built from \garnet's
  own site coordinates;
\item an XML material file with one mix per hexagon and plane: 55 hexagons
  $\times$ 10 planes of 20~cm;
\item the kinetics in a \code{PREC\_Filename} file.
\end{itemize}
The XML writer round-trips through \garnet's reader exactly. The only
exception is $D$, which passes through $\Sigma_{tr} = 1/(3D)$ as in \ffn{}
and so loses one ulp.

\paragraph{HP-MR transient: external reactivity insertion.} From the critical
initial state, every fissile $\nu\Sigma_f$ is multiplied by
$1 + \varepsilon(t)$. $\varepsilon$ rises linearly from 0 to
$\varepsilon_1 = \rho/(1-\rho)$ over $[0, 2]$~s, with
$\rho = 0.5\,\beta = 340.1$~pcm. Scaling all production scales $k$, so the
final state is exactly $+0.5$\,\$ statically on any mesh. $\Sigma_f$, the
power weight, is unchanged.

The kinetics are the library's 11 group speeds and six precursor families
($\beta = 680.2$~pcm), with the delayed spectrum set equal to $\chi$ in both
codes, and there is no feedback. \ffn{} runs this through
\code{Type\_Perturbation = Read\_XML\_File}, which interpolates every cross
section linearly from the initial to a final XML over $[0, T_\mathrm{end}]$;
here the two files differ only in \code{NuSigF}. \garnet's \code{problem\_at}
applies the same linear scaling.

\paragraph{Three \ffn{} transient findings.} All three were found with a null
transient (final XML equal to the initial one), which must keep $P = 1$.
\begin{enumerate}
\item With the kinetics in a \code{PREC\_Filename} file (the format of \ffn's
  C5G7-TD example), an XML material keeps a zero prompt spectrum. \ffn{}
  stores the prompt and delayed spectra of an XML material only when it has
  \code{ChiD}, and the unperturbed transient collapsed. The kinetics are
  therefore written into the XML.
\item The default Implicit-Exponential scheme takes the decay constants and
  delayed spectra of material~0 for the whole core
  (\code{get\_delayed\_decay\_constant(0, p)}). Our mix~0 is a non-fissile
  axial-reflector cell, so every non-fissile mix carries the global $\lambda$
  and the fissile $\chi$ with $\beta = 0$; all fissile mixes share one set
  of delayed data. This changes nothing physically.
\item Even so, Implicit-Exponential does not hold the null transient: $P$
  drops by 3.9\% in a 0.1~s step and by 0.2\% in a 0.01~s step, and keeps
  drifting. Both semi-implicit schemes (Semi-Implicit-Euler and
  Semi-Implicit-Exponential) hold $P = 1.0000$. The comparison uses those two.
\end{enumerate}

\paragraph{Why the drums are withdrawn.}\label{sec:blackdrum} The model was
first set up with each drum hexagon homogenized by area: the drum's $90^\circ$
B$_4$C arc covers $f_0 = 3.2$\% of the hexagon, with $\Sigma_t$ mixed and
$D = 1/(3\Sigma_t)$ (\code{absorber="smeared"}). The result is a hexagon that
is black to thermal neutrons: $\Sigma_a \approx 10$~cm$^{-1}$,
$D \approx 0.03$~cm, a diffusion length of about 0.5~mm. Neither code
converges on it at any practical mesh:
\begin{itemize}
\item \ffn{} FE1 to FE3 levels off near $k = 1.0347$;
\item \garnet{} came down from 1.0553 ($r=2$) through 1.0494 ($r=4$) to
  1.0466 ($r=6$, 7.8 million unknowns), roughly first order and still
  1200~pcm above.
\end{itemize}
With the drums made pure beryllium, the two codes agree to a few pcm at the
same meshes. That confirms the geometry, boundary conditions and cross-section
semantics. Such a drum would be a benchmark of resolving an unphysical
sub-millimetre boundary layer, not of discretization cost.

The same smearing rules out a drum-rotation transient: a linear ramp of the
smeared B$_4$C fraction inserts almost all of its reactivity in its last few
percent, because the hexagon stays black until the fraction falls below about
10\% of nominal. The transient is therefore an external reactivity insertion.
